\chapter*{Abstract}

\par Negai Marin, \textit{``Old Bugs, New Protocol: A Root-Cause Analysis of Security Failures in the Model Context Protocol Ecosystem Under Rapid Adoption''}, research article, Technical University of Moldova, Faculty of Computers, Informatics and Microelectronics, Chișinău, 2026. The article has an introduction, four chapters, conclusions, a bibliography of 30 sources and one appendix. It contains 5 figures and 5 tables.

\vspace{0.3cm}
\par \noindent \textbf{Keywords:} Model Context Protocol, AI agents, software security, prompt injection, vulnerability analysis.
\vspace{0.3cm}

\par The Model Context Protocol (MCP) was released at the end of 2024 as a common way to connect language-model assistants to files, databases and online services. Within a year it was supported by most large AI vendors, and during the same period security researchers published dozens of reports about vulnerable servers, compromised packages and attacks that hide instructions inside tool descriptions. These reports are usually discussed as a new ``AI security'' problem. The question behind this article is whether that label is accurate, or whether most of the failures are ordinary software defects that appeared because the ecosystem grew faster than anyone could review it.

\par The objective was to find where MCP security failures come from. I collected \nRecords{} publicly documented weaknesses and incidents disclosed between \firstDisclosure{} and \lastDisclosure{}, mainly from an open vulnerability catalogue and partly from the original research disclosures. Every record was read in its original source and coded on two axes. The first axis is the nature of the flaw: classic (it would exist even without a language model), AI-native (it works only because a model follows text it should not trust) or hybrid. The second axis is the stage where the flaw was introduced: protocol design, implementation, integration design, supply chain or deployment. All counts and charts were produced by a small Python script from the coded dataset, which is published together with the article.

\par The results show that \pClassic\% of the records are classic weaknesses such as command injection, path traversal, server-side request forgery and missing authentication, and half of all records are coding mistakes in individual servers, clients or tools. AI-native weaknesses make up \pNative\% of the sample and come almost entirely from the design of the protocol and of the integrations built on it. None of them received a CVE identifier, which means that tools based on vulnerability databases cannot see this part of the risk. The article ends with a set of controls for developers, operators and public institutions, and compares the findings with the October 2026 breach of Moldovan e-government services, where a comparable trust error led to the exposure of personal data of about 450,000 citizens.
